\section{Conclusion}

Under a protocol with a shared learning rate and checkpoint selection by validation balanced accuracy at 0.5, focal loss with class-balanced sampling appeared to help a vision transformer and hurt two CNNs in melanoma classification. The effect agreed across seeds, one architecture interaction excluded zero, and an external dataset pointed the same way for the CNNs. Under a revised protocol with a per-architecture learning rate, threshold-free checkpoint selection, and five seeds, the interaction disappeared. The correction lowered sensitivity by 1.7 to 4.0 points in all three architectures, with intervals that include zero, and did not detectably change the AUC. The ViT, undertrained under the original protocol, had the highest point-estimate AUC in domain and on HIBA.

The skin-lightness findings did not depend on the protocol. After adjustment for acquisition source and lesion contrast, lighter skin was associated with more false negatives in both CNNs and darker skin with more false positives in all three architectures. The darkest stratum had the lowest AUC, the highest false positive rate, the worst calibration, and the lowest Grad-CAM pointing-game accuracy in every architecture, although the per-run AUC differences were not significant and the stratum contained six melanomas. On HIBA, imbalance correction had no detectable effect, while strong augmentation improved external AUC in every seed and sensitivity in 14 of 15 runs.

Architecture comparisons under training interventions should tune each architecture, select checkpoints with a rule that does not depend on the score scale, and treat an isolated interaction as a hypothesis until it survives a change of protocol. The main requirement for a precise claim about Fitzpatrick V-VI skin remains dermoscopic data from those groups.
